% =====================================================================
% Figure — §V trajectory comparison: static constraint vs VA-ATACOM.
%
% Generated by
% experiments/diagnostics/F1_trajectory/plot_f1.py.
% A double-integrator (relative-degree-2) simulation with a single
% circular obstacle; a velocity-setpoint policy seeks the goal and the
% executed velocity is rate-limited at a_max (momentum), so the filter
% alone shapes the trajectory. Native canvas is 7.0 × 3.0 in.
%
% Usage in §V (qualitative):
%   % =====================================================================
% Figure — §V trajectory comparison: static constraint vs VA-ATACOM.
%
% Generated by
% experiments/diagnostics/F1_trajectory/plot_f1.py.
% A double-integrator (relative-degree-2) simulation with a single
% circular obstacle; a velocity-setpoint policy seeks the goal and the
% executed velocity is rate-limited at a_max (momentum), so the filter
% alone shapes the trajectory. Native canvas is 7.0 × 3.0 in.
%
% Usage in §V (qualitative):
%   % =====================================================================
% Figure — §V trajectory comparison: static constraint vs VA-ATACOM.
%
% Generated by
% experiments/diagnostics/F1_trajectory/plot_f1.py.
% A double-integrator (relative-degree-2) simulation with a single
% circular obstacle; a velocity-setpoint policy seeks the goal and the
% executed velocity is rate-limited at a_max (momentum), so the filter
% alone shapes the trajectory. Native canvas is 7.0 × 3.0 in.
%
% Usage in §V (qualitative):
%   % =====================================================================
% Figure — §V trajectory comparison: static constraint vs VA-ATACOM.
%
% Generated by
% experiments/diagnostics/F1_trajectory/plot_f1.py.
% A double-integrator (relative-degree-2) simulation with a single
% circular obstacle; a velocity-setpoint policy seeks the goal and the
% executed velocity is rate-limited at a_max (momentum), so the filter
% alone shapes the trajectory. Native canvas is 7.0 × 3.0 in.
%
% Usage in §V (qualitative):
%   \input{figures/fig_v_a_trajectory/fig_v_a_trajectory.tex}
%
% Cross-references: Prop. 1-2 (M1/M2), §IV-B Component 1, Table 2.
% =====================================================================
\begin{figure*}[t]
  \centering
  \includegraphics[width=0.95\textwidth]%
    {figures/fig_v_a_trajectory/fig_v_a_trajectory.pdf}
  \caption{%
    \textbf{Why velocity augmentation matters (double-integrator
    simulation, $\Delta t{=}0.1$\,s, $v_{\max}{=}1$\,m/s,
    $a_{\max}{=}0.9$\,m/s$^2$).} Agent starts at $(-1,0)$ with goal
    $(1,0)$ and an obstacle of radius $r{=}0.2$\,m just off the
    start--goal line; a velocity-setpoint policy seeks the goal and the
    executed velocity is rate-limited at $a_{\max}$ (momentum). Both
    panels use \emph{identical} dynamics and the \emph{same} tangent
    projection --- the only difference is whether the constraint
    includes the braking term $\|\boldsymbol{v}\|^2/(2a_{\max})$.
    \emph{Left, static constraint $\rho \ge r_\mathrm{base}$:} the
    filter does not engage until the boundary, by which point the
    committed inward velocity cannot be cancelled in one step; four
    consecutive steps (red) penetrate to $\rho_{\min}{=}0.076$\,m
    (collision) --- the M1/M2 momentum regime of Prop.~1--2.
    \emph{Right, VA-ATACOM:} the braking term caps the inward approach
    speed at the brakeable value $\sqrt{2a_{\max}(\rho-r_\mathrm{base})}$,
    so the agent decelerates early and arcs around at
    $\rho_{\min}{=}0.28$\,m before re-acquiring the goal --- the velocity
    augmentation isolated in the design study (Table~\ref{tab:design}).%
  }
  \label{fig:v_a_trajectory}
\end{figure*}

%
% Cross-references: Prop. 1-2 (M1/M2), §IV-B Component 1, Table 2.
% =====================================================================
\begin{figure*}[t]
  \centering
  \includegraphics[width=0.95\textwidth]%
    {figures/fig_v_a_trajectory/fig_v_a_trajectory.pdf}
  \caption{%
    \textbf{Why velocity augmentation matters (double-integrator
    simulation, $\Delta t{=}0.1$\,s, $v_{\max}{=}1$\,m/s,
    $a_{\max}{=}0.9$\,m/s$^2$).} Agent starts at $(-1,0)$ with goal
    $(1,0)$ and an obstacle of radius $r{=}0.2$\,m just off the
    start--goal line; a velocity-setpoint policy seeks the goal and the
    executed velocity is rate-limited at $a_{\max}$ (momentum). Both
    panels use \emph{identical} dynamics and the \emph{same} tangent
    projection --- the only difference is whether the constraint
    includes the braking term $\|\boldsymbol{v}\|^2/(2a_{\max})$.
    \emph{Left, static constraint $\rho \ge r_\mathrm{base}$:} the
    filter does not engage until the boundary, by which point the
    committed inward velocity cannot be cancelled in one step; four
    consecutive steps (red) penetrate to $\rho_{\min}{=}0.076$\,m
    (collision) --- the M1/M2 momentum regime of Prop.~1--2.
    \emph{Right, VA-ATACOM:} the braking term caps the inward approach
    speed at the brakeable value $\sqrt{2a_{\max}(\rho-r_\mathrm{base})}$,
    so the agent decelerates early and arcs around at
    $\rho_{\min}{=}0.28$\,m before re-acquiring the goal --- the velocity
    augmentation isolated in the design study (Table~\ref{tab:design}).%
  }
  \label{fig:v_a_trajectory}
\end{figure*}

%
% Cross-references: Prop. 1-2 (M1/M2), §IV-B Component 1, Table 2.
% =====================================================================
\begin{figure*}[t]
  \centering
  \includegraphics[width=0.95\textwidth]%
    {figures/fig_v_a_trajectory/fig_v_a_trajectory.pdf}
  \caption{%
    \textbf{Why velocity augmentation matters (double-integrator
    simulation, $\Delta t{=}0.1$\,s, $v_{\max}{=}1$\,m/s,
    $a_{\max}{=}0.9$\,m/s$^2$).} Agent starts at $(-1,0)$ with goal
    $(1,0)$ and an obstacle of radius $r{=}0.2$\,m just off the
    start--goal line; a velocity-setpoint policy seeks the goal and the
    executed velocity is rate-limited at $a_{\max}$ (momentum). Both
    panels use \emph{identical} dynamics and the \emph{same} tangent
    projection --- the only difference is whether the constraint
    includes the braking term $\|\boldsymbol{v}\|^2/(2a_{\max})$.
    \emph{Left, static constraint $\rho \ge r_\mathrm{base}$:} the
    filter does not engage until the boundary, by which point the
    committed inward velocity cannot be cancelled in one step; four
    consecutive steps (red) penetrate to $\rho_{\min}{=}0.076$\,m
    (collision) --- the M1/M2 momentum regime of Prop.~1--2.
    \emph{Right, VA-ATACOM:} the braking term caps the inward approach
    speed at the brakeable value $\sqrt{2a_{\max}(\rho-r_\mathrm{base})}$,
    so the agent decelerates early and arcs around at
    $\rho_{\min}{=}0.28$\,m before re-acquiring the goal --- the velocity
    augmentation isolated in the design study (Table~\ref{tab:design}).%
  }
  \label{fig:v_a_trajectory}
\end{figure*}

%
% Cross-references: Prop. 1-2 (M1/M2), §IV-B Component 1, Table 2.
% =====================================================================
\begin{figure*}[t]
  \centering
  \includegraphics[width=0.95\textwidth]%
    {figures/fig_v_a_trajectory/fig_v_a_trajectory.pdf}
  \caption{%
    \textbf{Why velocity augmentation matters (double-integrator
    simulation, $\Delta t{=}0.1$\,s, $v_{\max}{=}1$\,m/s,
    $a_{\max}{=}0.9$\,m/s$^2$).} Agent starts at $(-1,0)$ with goal
    $(1,0)$ and an obstacle of radius $r{=}0.2$\,m just off the
    start--goal line; a velocity-setpoint policy seeks the goal and the
    executed velocity is rate-limited at $a_{\max}$ (momentum). Both
    panels use \emph{identical} dynamics and the \emph{same} tangent
    projection --- the only difference is whether the constraint
    includes the braking term $\|\boldsymbol{v}\|^2/(2a_{\max})$.
    \emph{Left, static constraint $\rho \ge r_\mathrm{base}$:} the
    filter does not engage until the boundary, by which point the
    committed inward velocity cannot be cancelled in one step; four
    consecutive steps (red) penetrate to $\rho_{\min}{=}0.076$\,m
    (collision) --- the M1/M2 momentum regime of Prop.~1--2.
    \emph{Right, VA-ATACOM:} the braking term caps the inward approach
    speed at the brakeable value $\sqrt{2a_{\max}(\rho-r_\mathrm{base})}$,
    so the agent decelerates early and arcs around at
    $\rho_{\min}{=}0.28$\,m before re-acquiring the goal --- the velocity
    augmentation isolated in the design study (Table~\ref{tab:design}).%
  }
  \label{fig:v_a_trajectory}
\end{figure*}
